%% Abstract.
Flux is a decentralized cloud of \num{6489} collateralised nodes running \num{1195} deployed
applications, on a chain that replaced proof-of-work with collateral-gated sortition at height
\num{2020000}. This paper does two things. It describes the deployed system --- consensus, node
tiering and benchmarking, monetary policy, orchestration and placement, addressing, attested
execution, and the Zcash-derived shielded layer --- at a depth from which the behaviour could be
reimplemented or audited, citing source by file and line --- with the caveat, stated in full in
\S\ref{sec:limitations}, that a minority of the components cited are private repositories a reader
cannot open, and that one class of finding in those is deliberately stated without coordinates. It then specifies the target architecture
at the same depth: revenue-funded operator compensation replacing issuance, a redesigned bridge,
a content-reporting quorum, application specification v9, a bonded attestation network,
post-quantum migration, and the machine-facing interfaces an autonomous tenant needs to provision,
pay for and release capacity without a human. It also argues one positioning claim from the
substrate rather than from intent: the interconnect requirement that centralises frontier-model
inference does not apply to models that fit on a single device, so a dispersed fleet of
collateralised commodity machines is a structurally appropriate home for small task-specialised
models --- once accelerators can be subdivided, which is the one thing the current allocation
model cannot do.

The discipline that keeps the second half honest is that every undecided element is named rather
than filled: each mechanism is given as state, transitions, invariants and failure handling, with
open parameters marked as decisions and open constructions marked as open. Two axes of tagging let a
reader tell, at any sentence, whether they are reading about code that runs or a design that does
not.

Several findings are not what the roadmap or the design intent would predict, and the paper reports
them as measured. A fee pool with a per-block drip is necessary because deterministic payment
rotation otherwise admits a timing attack recapturing most of an operator's own fee; the same pool
design pays every node regardless of what it hosts, so it supplies no gradient rewarding good
hosting, and what compels performance is the eligibility gate --- which the deployed benchmark and
attestation path does not bind to a machine. Hardening that gate is urgent independently of
revenue-sharing rather than because of it: at activation the gate would already be admitting
\SI{95.7}{\percent} of operator income as issuance, and the new mechanism would add the other
\SI{4.3}{\percent}. The distribution of that revenue would be a consensus rule;
\emph{which} revenue is subject to it, and at what price, would remain policy, and the paper marks
that boundary wherever the decentralization claim rests on it. Collateral-weighted moderation is Sybil-resistant by construction and, at the
measured ownership distribution, not censorship-resistant: four owner identities reach the proposed
removal threshold, and no weighting or remedy in the design space raises that number. The shielded
pools have been closed to new entrants by consensus since 2021 and hold \SI{0.0225}{\percent} of
supply; they are to be retired rather than reopened, leaving \Flux{} settling transparently and
treating payment privacy as an unsolved problem rather than a delivered feature. A header format chosen for consensus reasons forecloses light-client
bridging. Of twenty-two mechanisms assessed, six are extensions of what exists, eleven are
engineering, and four are open research problems
--- which the paper names rather than absorbs.
